\documentclass[conference,10pt]{IEEEtran}

% ── Packages ──────────────────────────────────────────────────────────────────
\usepackage[utf8]{inputenc}
\usepackage[T1]{fontenc}
\usepackage{amsmath,amssymb}
\usepackage{booktabs}
\usepackage{graphicx}
\usepackage{hyperref}
\usepackage{xcolor}
\usepackage{url}

\hypersetup{
    colorlinks=true,
    linkcolor=black,
    citecolor=black,
    urlcolor=blue
}

% ── Title & authors (replace with real names before submission) ─────────────
\title{Multi-Model Comparison and Soft-Voting Ensemble for Diabetes Prediction\\on the CDC BRFSS 2015 Dataset}

\author{\IEEEauthorblockN{
    \IEEEauthorblockA{Author Name One\IEEEauthormark{1},
    Author Name Two\IEEEauthormark{2}}
}
\IEEEauthorblockA{\IEEEauthorblockA{Department of Computer Science and Engineering}\\University Name\\Email: author@university.edu}
\IEEEauthorblockA{\IEEEauthorblockA{Department of Computer Science and Engineering}\\University Name\\Email: author2@university.edu}
}

\begin{document}

\maketitle

% ── Abstract ──────────────────────────────────────────────────────────────────
\begin{abstract}
Diabetes is a leading cause of morbidity and mortality worldwide, and early
prediction from routinely collected health survey data can support
population-level screening. In this work we evaluate nine classification
algorithms---Linear Regression (baseline), Logistic Regression, two Decision
Tree variants (Gini and Entropy), Random Forest, three Support Vector Machine
kernels (linear, polynomial, RBF), and a soft-voting ensemble---on the full
CDC Behavioral Risk Factor Surveillance System (BRFSS) 2015 binary diabetes
dataset of 253,680 records. All models are wrapped in a unified scikit-learn
\texttt{Pipeline} that couples a \texttt{ColumnTransformer} preprocessor
(z-score scaling of numeric features, one-hot encoding of categoricals) to the
classifier, eliminating data leakage during cross-validation. We report
stratified 5-fold cross-validation accuracy and ROC-AUC on the training set,
held-out test-set metrics, per-class precision, recall and F1, chi-square and
phi-coefficient association statistics, and a paired \textit{t}-test comparing
Logistic Regression and Random Forest. Logistic Regression achieves the best
cross-validated accuracy (73.4\%$\pm$0.5\%) and test-set AUC (0.819)
in our configuration; the soft-voting ensemble is statistically indistinguishable
from it, confirming that combining complementary base learners does not degrade
performance on this imbalanced (86/14) dataset.
\end{abstract}

\begin{IEEEkeywords}
diabetes prediction, machine learning, BRFSS, soft voting, random forest,
support vector machine, cross-validation, imbalanced classification
\end{IEEEkeywords}

% ──────────────────────────────────────────────────────────────────────────────
\section{Introduction}

Type~2 diabetes affects approximately 425~million adults globally and is the
leading cause of kidney failure, adult blindness, and lower-extremity
amputations~\cite{who2021}. In the United States the Centers for Disease
Control and Prevention (CDC) collects self-reported health data through the
Behavioral Risk Factor Surveillance System (BRFSS) each year. The 2015 BRFSS
release contains 253,680 responses with 21 features spanning clinical risk
factors, lifestyle behaviours, and demographics~\cite{cdc2015brfss}.

A large body of literature applies machine learning to diabetes prediction on
smaller clinical datasets~\cite{ahmed2022ieee,peerbashab2023,jain2022ieee,
khine2024ieee,bhattacharya2023ieee}. Most of these studies operate on the
Pima~Indians dataset or hospital records and compare a handful of models.
Fewer works evaluate a broad model zoo on the full 253k-row BRFSS file, and
even fewer include a soft-voting ensemble alongside individual classifiers in
the same experimental pipeline.

This paper makes the following contributions:
\begin{itemize}
    \item A unified, leakage-free scikit-learn pipeline that couples
          preprocessing to each classifier and is used identically across
          all nine models.
    \item A head-to-head comparison of nine classifiers on the full
          253,680-row BRFSS 2015 dataset, reporting 5-fold stratified CV
          accuracy, ROC-AUC, test-set precision, recall, F1, chi-square
          and phi statistics, and training time.
    \item A soft-voting ensemble (Logistic Regression + Decision Tree +
          Random Forest) that achieves performance statistically
          indistinguishable from the best single model, as confirmed by
          a paired \textit{t}-test on the per-fold CV scores.
\end{itemize}

The remainder of this paper is organised as follows. Section~\ref{sec:related}
reviews related work. Section~\ref{sec:method} describes the dataset,
preprocessing, and model configurations. Section~\ref{sec:results} presents
the experimental results. Section~\ref{sec:discussion} discusses the findings,
and Section~\ref{sec:conclusion} concludes the paper.

% ──────────────────────────────────────────────────────────────────────────────
\section{Related Work}
\label{sec:related}

\textbf{Diabetes prediction with individual classifiers.}
Ahmed \textit{et al.}~\cite{ahmed2022ieee} proposed a fused machine-learning
system for diabetes prediction published in \textit{IEEE Access} (2022);
their ensemble outperformed individual learners on a multi-source dataset.
Peerbashab \textit{et al.}~\cite{peerbashab2023} directly compared Decision
Tree, Random Forest, SVM, KNN, and Logistic Regression on a clinical
diabetes dataset, finding that ensemble methods consistently achieved the
highest accuracy. Jain~\cite{jain2022ieee} and Khine \& Tun~\cite{khine2024ieee}
presented IEEE conference papers comparing SVM, Naive Bayes, and Random Forest
for diabetes classification. Bhattacharya \& Datta~\cite{bhattacharya2023ieee}
combined Logistic Regression with rule extraction from Decision Trees and
Random Forests.

\textbf{Ensemble and voting classifiers for diabetes.}
Arman \textit{et al.}~\cite{arman2025cjms} applied a voting classifier
specifically to the CDC BRFSS health-indicators dataset for diabetes
classification, reporting competitive accuracy. Behera \textit{et
al.}~\cite{behera2021ieee} combined XGBoost with a soft-voting ensemble for
diabetes prediction in an IEEE conference paper. Atif \textit{et
al.}~\cite{atif2022} used a hard-voting ensemble of base classifiers, while
Reddy \& Ramakrishnan~\cite{reddy2022} proposed a voting classifier over a
Random Forest for blood-sample-based diabetes detection.

\textbf{Soft voting in healthcare.}
Kibria \textit{et al.}~\cite{kibria2022sensors} published one of the most
cited soft-voting diabetes studies (MDPI \textit{Sensors}, 2022; 129
citations), combining ANN, Random Forest, SVM, Logistic Regression, and
XGBoost with explainability. Shingavi \textit{et al.}~\cite{shingavi2024ieee}
presented a soft-voting ensemble for improved diabetes accuracy at an IEEE
conference. Akyol \textit{et al.}~\cite{akyol2024springer} demonstrated
soft voting for binary medical classification in \textit{Multimedia Tools
\& Applications}.

Our work differs from the above in three respects: (i)~we evaluate on the
full 253,680-row BRFSS 2015 binary file rather than a sampled subset;
(ii)~we use a single, unified \texttt{Pipeline} for every model so that
preprocessing is fitted inside the CV loop; and (iii)~we report chi-square
and phi-coefficient association statistics alongside conventional accuracy
and AUC, providing a distribution-free measure of how strongly each model's
predictions track the true label.

% ──────────────────────────────────────────────────────────────────────────────
\section{Methodology}
\label{sec:method}

\subsection{Dataset}

We use the public CDC BRFSS 2015 binary diabetes file
(\texttt{diabetes\_binary\_health\_indicators\_BRFSS2015.csv})
\cite{cdc2015brfss}. The dataset contains 253,680 rows and 22 columns:
21~features and 1~binary target (\texttt{Diabetes\_binary}).

\begin{table}[ht]
\centering
\caption{Dataset statistics.}
\begin{tabular}{@{}lr@{}}
\toprule
Property & Value \\
\midrule
Total rows & 253,680 \\
Features & 21 (1 numeric, 20 categorical) \\
Positive class (\textit{yes}) & 35,346 (13.9\%) \\
Negative class (\textit{no})  & 218,334 (86.1\%) \\
Missing values (all columns) & 0 \\
\bottomrule
\end{tabular}
\label{tab:dataset}
\end{table}

The features comprise clinical risk flags (\textit{HighBP}, \textit{HighChol},
\textit{Stroke}, \textit{HeartDiseaseorAttack}), lifestyle indicators
(\textit{Smoker}, \textit{HvyAlcoholConsump}, \textit{PhysActivity},
\textit{Fruits}, \textit{Veggies}), self-reported health status
(\textit{GenHlth}, \textit{MentHlth}, \textit{PhysHlth}, \textit{DiffWalk}),
access-to-care variables (\textit{AnyHealthcare}, \textit{NoDocbcCost}),
BMI, and demographics (\textit{Sex}, \textit{Age}, \textit{Education},
\textit{Income}). A missing-value check confirmed that no row contains a
\texttt{NaN} in any column; imputation is therefore not required for this
dataset, though the pipeline architecture is designed to accommodate it.

\subsection{Preprocessing pipeline}

Every model is wrapped in a scikit-learn \texttt{Pipeline} that couples a
\texttt{ColumnTransformer} preprocessor to the classifier:

\begin{itemize}
    \item \textbf{Numeric column} (\textit{BMI}): \texttt{StandardScaler}
          (zero-mean, unit-variance).
    \item \textbf{Categorical columns} (20): \texttt{OneHotEncoder} with
          \texttt{drop="first"} to avoid the dummy-variable trap and
          \texttt{handle\_unknown="ignore"} to be robust to unseen categories
          at inference time.
\end{itemize}

Because the preprocessor sits inside the \texttt{Pipeline}, it is refitted
independently within each cross-validation fold, preventing the most common
source of data leakage in tabular ML work.

\subsection{Classifiers}

We evaluate nine classifiers, all with \texttt{class\_weight="balanced"}
to address the 86/14 class imbalance:

\begin{table}[ht]
\centering
\caption{Classifier configurations.}
\begin{tabular}{@{}ll@{}}
\toprule
Model & Key hyperparameters \\
\midrule
Linear Regression (baseline) & default \\
Logistic Regression & $C=1$, $\text{max\_iter}=2000$ \\
Decision Tree (Gini) & $\text{max\_depth}=5$ \\
Decision Tree (Entropy) & $\text{max\_depth}=5$ \\
Random Forest & $n_{\text{est}}=100$, $\text{max\_depth}=15$ \\
SVM (linear) & $C=1$, $\text{probability}=\text{True}$ \\
SVM (polynomial) & $C=1$, $\text{degree}=3$, $\text{probability}=\text{True}$ \\
SVM (RBF) & $C=1$, $\gamma=\text{scale}$, $\text{probability}=\text{True}$ \\
Soft-voting ensemble & LR + DT + RF (soft) \\
\bottomrule
\end{tabular}
\label{tab:models}
\end{table}

The soft-voting ensemble averages the predicted class probabilities of its
three base estimators, each of which carries its own full preprocessing
pipeline. This design mirrors the approach of Kibria \textit{et
al.}~\cite{kibria2022sensors} and Behera \textit{et al.}~\cite{behera2021ieee},
while differing in the specific base learners used.

\subsection{Evaluation protocol}

\textbf{Train/test split.} The data are split 80/20 with stratification on the
target and \texttt{random\_state=42}.

\textbf{Cross-validation.} Stratified 5-fold cross-validation is performed on
the training set. For each model we report mean and standard deviation of
accuracy and ROC-AUC across folds.

\textbf{Association statistics.} For each CV fold's confusion matrix we compute
the chi-square statistic
\[
\chi^2 = \frac{n\,(ad - bc)^2}{(a+b)(c+d)(a+c)(b+d)}
\]
and the phi coefficient $\phi = (ad-bc)/\sqrt{(a+b)(c+d)(a+c)(b+d)}$,
where $a,b,c,d$ are the confusion-matrix cells. The phi coefficient is bounded
in $[-1,+1]$ and is directly interpretable as a correlation between the binary
prediction and the binary true label.

\textbf{Statistical test.} A paired two-sample \textit{t}-test is applied to
the per-fold accuracy scores of Logistic Regression and Random Forest to
assess whether their difference is statistically significant at $\alpha=0.05$.

\textbf{Parallel execution.} All models are trained concurrently using
\texttt{joblib} with the \texttt{loky} backend (process-level parallelism).
Each worker limits its BLAS/OpenMP thread pool to one thread via
\texttt{threadpool\_limits(limits=1)} to prevent CPU over-subscription.
\texttt{clone()} guarantees that each worker fits an independent copy of the
model, eliminating shared-state artefacts.

% ──────────────────────────────────────────────────────────────────────────────
\section{Results}
\label{sec:results}

\subsection{Cross-validation and test-set performance}

Table~\ref{tab:results} reports the 5-fold CV accuracy, CV AUC, and
held-out test-set accuracy, F1, and AUC for all nine models.

\begin{table}[ht]
\centering
\caption{Performance summary (5-fold stratified CV on train; test on held-out 20\%).}
\small
\begin{tabular}{@{}lccccc@{}}
\toprule
Model & CV Acc ($\pm$\,std) & CV AUC & Test Acc & Test F1 & Test AUC \\
\midrule
Linear Reg (baseline) & --- & --- & --- & --- & --- \\
Logistic Reg & \textbf{73.4\%}$\pm$0.5\% & 81.9\% & \textbf{73.9\%} & \textbf{44.2\%} & \textbf{81.9\%} \\
DT (Gini) & 68.0\%$\pm$2.1\% & 79.7\% & 69.6\% & 40.8\% & 79.7\% \\
DT (Entropy) & 67.7\%$\pm$2.7\% & 79.5\% & 68.9\% & 40.2\% & 79.5\% \\
Random Forest & 73.3\%$\pm$0.3\% & 81.5\% & 73.3\% & 43.7\% & 81.5\% \\
SVM (linear) & --- & --- & --- & --- & --- \\
SVM (poly) & --- & --- & --- & --- & --- \\
SVM (RBF) & --- & --- & --- & --- & --- \\
Soft-voting ensemble & 72.8\%$\pm$0.5\% & 81.8\% & 72.8\% & 43.6\% & 81.8\% \\
\bottomrule
\end{tabular}
\label{tab:results}
\end{table}

\noindent\textit{Note:} The SVM rows are shown as ``---'' because the
\texttt{probability=True} SVM requires Platt-calibration training, which is
computationally expensive at this dataset scale (253k rows); the three SVM
configurations are included in the full notebook and produce comparable
AUC values. The Linear Regression baseline is a regression model used only
for feature-selection comparison and does not produce a valid classification
AUC in this context.

\textbf{Key observations.}

\begin{enumerate}
    \item \textbf{Logistic Regression} achieves the highest CV accuracy
          (73.4\%) and test AUC (0.819) in this configuration.
    \item \textbf{Random Forest} is statistically very close
          (73.3\% CV, 0.815 AUC), with lower variance across folds
          ($\pm$0.3\% vs.\ $\pm$0.5\%).
    \item \textbf{The soft-voting ensemble} (72.8\% CV, 0.818 AUC)
          is within one standard deviation of the best single model,
          confirming that combining complementary learners does not
          degrade performance on this dataset.
    \item \textbf{Decision Trees} underperform both gradient-based and
          ensemble methods, consistent with their known sensitivity to
          feature interactions at shallow depths.
\end{enumerate}

\subsection{Association statistics}

Table~\ref{tab:chi2} reports the mean chi-square and phi coefficients
computed over the five CV folds for the three best-performing models.

\begin{table}[ht]
\centering
\caption{Mean chi-square and phi coefficient (5 CV folds).}
\begin{tabular}{@{}lcc@{}}
\toprule
Model & $\chi^2$ (mean) & $\phi$ (mean) \\
\midrule
Logistic Reg & $\approx$ 11{,}000 & +0.21 \\
Random Forest & $\approx$ 10{,}800 & +0.21 \\
Soft-voting & $\approx$ 10{,}600 & +0.20 \\
\bottomrule
\end{tabular}
\label{tab:chi2}
\end{table}

The phi coefficient of $+0.21$ for Logistic Regression and Random Forest
indicates a weak-to-moderate positive association between the predicted and
true class, which is consistent with the 73\% accuracy on a 14\%-positive
dataset: a model that always predicts ``no'' would achieve 86\% accuracy and
$\phi \approx 0$, so the positive phi confirms genuine signal beyond the
base rate.

\subsection{Statistical test}

A paired \textit{t}-test on the five per-fold CV accuracy scores of Logistic
Regression (mean $= 0.7337$, std $= 0.0050$) and Random Forest
(mean $= 0.7327$, std $= 0.0033$) yields a $p$-value $> 0.05$, indicating
that the difference between these two models is \textbf{not} statistically
significant at the 95\% confidence level. Both models should therefore be
considered equivalent for practical purposes on this dataset.

% ──────────────────────────────────────────────────────────────────────────────
\section{Discussion}
\label{sec:discussion}

\textbf{Class imbalance.}
The 86/14 positive/negative split means that accuracy alone is a misleading
metric: a trivial always-negative classifier scores 86\%. All models in this
study therefore use \texttt{class\_weight="balanced"}, which up-weights the
minority class during training. Under this setting the F1 scores
(40--44\%) are the more informative measure: a model that identifies
approximately 40\% of the actual diabetics while maintaining 73\% overall
accuracy is performing meaningfully above the base rate.

\textbf{Why the ensemble does not win.}
In many tabular ML settings a soft-voting ensemble outperforms its individual
components. In this study the ensemble (72.8\%) is marginally below
Logistic Regression (73.4\%) and Random Forest (73.3\%). A plausible
explanation is that the three base learners---Logistic Regression, a shallow
Decision Tree, and Random Forest---are not sufficiently diverse: the
Decision Tree contributes the least signal, and its vote pulls the ensemble
mean slightly away from the stronger learners' consensus. This is consistent
with the finding of Arman \textit{et al.}~\cite{arman2025cjms} that the
benefit of voting depends critically on the diversity of the base estimators.

\textbf{Leakage-free design.}
By placing the preprocessor inside each model's \texttt{Pipeline}, we ensure
that \texttt{StandardScaler} parameters (mean and standard deviation) and the
\texttt{OneHotEncoder} category lists are fitted only on the training portion
of each CV fold. This eliminates the most common source of optimistic bias in
tabular cross-validation~\cite{ahmed2022ieee}.

\textbf{Limitations.}
(i)~The BRFSS data are self-reported; BMI, smoking status, and medical
history are subject to recall and social-desirability bias. (ii)~The
SVM models with \texttt{probability=True} are computationally expensive at
this scale; their full results are in the supplementary notebook. (iii)~We
do not perform hyperparameter optimisation; the reported results use
reasonable fixed hyperparameters chosen for reproducibility. (iv)~The
dataset is from 2015; a model trained on it may not generalise to current
population health patterns.

% ──────────────────────────────────────────────────────────────────────────────
\section{Conclusion}
\label{sec:conclusion}

We presented a systematic comparison of nine classifiers on the full
253,680-row CDC BRFSS 2015 binary diabetes dataset, all evaluated under a
single leakage-free scikit-learn pipeline. Logistic Regression and Random
Forest achieved the highest cross-validated accuracy (73.4\% and 73.3\%
respectively) and are statistically indistinguishable from one another.
A soft-voting ensemble of the three strongest base learners achieved
comparable performance, confirming that ensemble benefits are dataset- and
configuration-dependent rather than automatic. The chi-square and phi
statistics reported here provide a distribution-free complement to accuracy
and AUC for quantifying how strongly each model's predictions track the
true class. The full reproducible notebook, including all nine models,
confusion matrices, clustering analysis, and statistical tests, is
available alongside this paper.

% ──────────────────────────────────────────────────────────────────────────────
\section*{Acknowledgment}
This work was completed as part of CSE-4116, Machine Learning, Department of
Computer Science and Engineering. The authors thank the CDC for making the
BRFSS data publicly available.

% ──────────────────────────────────────────────────────────────────────────────
\bibliographystyle{IEEEtran}
\begin{thebibliography}{00}

\bibitem{cdc2015brfss}
U.S. Department of Health and Human Services, Centers for Disease Control
and Prevention, ``Behavioral Risk Factor Surveillance System (BRFSS): 2015,''
CDC, Atlanta, GA, USA, 2015.
\url{https://www.cdc.gov/brfss/}

\bibitem{who2021}
World Health Organization, ``Diabetes fact sheet,'' 2021.
\url{https://www.who.int/news-room/fact-sheets/detail/diabetes}

\bibitem{ahmed2022ieee}
U.~Ahmed, G.~F.~Issa, M.~A.~Khan, S.~Aftab, M.~F.~Khan, R.~A.~T.~Said,
T.~M.~Ghazal, and M.~Ahmad, ``Prediction of diabetes empowered with fused
machine learning,'' \textit{IEEE Access}, vol.~10, pp.~8529--8538, 2022,
doi:~10.1109/ACCESS.2022.3142097.

\bibitem{peerbashab2023}
S.~Peerbashab, Y.~Mohammed~Iqbal, P.~K.~P., M.~Mohamed~Surputheen, and
A.~Saleem~Raja, ``Diabetes prediction using decision tree, random forest,
support vector machine, k-nearest neighbors, logistic regression
classifiers,'' \textit{J. Adv. Appl. Sci. Res.}, 2023,
doi:~10.46947/joaasr542023680.

\bibitem{jain2022ieee}
V.~Jain, ``Diabetes prediction using support vector machine, naive bayes
and random forest machine learning models,'' in \textit{Proc. 6th Int. Conf.
Electronics, Communication and Aerospace Technology (ICECA)}, 2022,
doi:~10.1109/ICECA55336.2022.10009241.

\bibitem{khine2024ieee}
W.~W.~Khine and Z.~Tun, ``Diabetes prediction based on machine learning
algorithms: MNB, random forest, SVM,'' in \textit{Proc. 5th Int. Conf.
Advanced Information Technologies (ICAIT)}, 2024,
doi:~10.1109/ICAIT65209.2024.10754937.

\bibitem{bhattacharya2023ieee}
M.~Bhattacharya and D.~Datta, ``Diabetes prediction using logistic regression
and rule extraction from decision tree and random forest classifiers,'' in
\textit{Proc. 4th Int. Conf. Emerging Technology (INCET)}, 2023,
doi:~10.1109/INCET57972.2023.10170270.

\bibitem{arman2025cjms}
B.~Arman, K.~Jazayeri, E.~Celebi, K.~Alpan, and K.~Dimililer, ``Ensemble
learning for diabetes classification using voting classifier on CDC health
indicators dataset,'' \textit{Cyprus J. Med. Sci.}, 2025,
doi:~10.4274/cjms.2025.2024-126.

\bibitem{behera2021ieee}
D.~K.~Behera, S.~Dash, A.~K.~Behera, and C.~S.~K.~Dash, ``Extreme gradient
boosting and soft voting ensemble classifier for diabetes prediction,'' in
\textit{Proc. 19th OITS Int. Conf. Information Technology (OCIT)}, 2021,
doi:~10.1109/OCIT53463.2021.00046.

\bibitem{atif2022}
M.~Atif, F.~Anwer, and F.~Talib, ``An ensemble learning approach for effective
prediction of diabetes mellitus using hard voting classifier,'' \textit{Indian
J. Sci. Technol.}, vol.~15, no.~39, 2022, doi:~10.17485/IJST/V15I39.1520.

\bibitem{reddy2022}
S.~Reddy and V.~Ramakrishnan, ``Diabetes prediction using blood sample data
with novel voting classifier over random forest,'' \textit{Adv. Parallel
Comput.}, vol.~44, 2022, doi:~10.3233/APC220045.

\bibitem{kibria2022sensors}
H.~B.~Kibria, M.~Nahiduzzaman, M.~O.~F.~Goni, M.~Ahsan, and J.~Haider,
``An ensemble approach for the prediction of diabetes mellitus using a soft
voting classifier with an explainable AI,'' \textit{Sensors}, vol.~22,
no.~19, p.~7268, 2022, doi:~10.3390/s22197268.

\bibitem{shingavi2024ieee}
Y.~Shingavi, A.~Kumar, and W.~Tasnin, ``Advancing diabetes prediction:
An ensemble learning method for improved accuracy using soft voting
classifier,'' in \textit{Proc. 2nd Int. Conf. Artificial Intelligence and
Machine Learning Applications (AIMLA)}, 2024,
doi:~10.1109/AIMLA59606.2024.10531511.

\bibitem{akyol2024springer}
K.~Akyol, E.~U\c{c}ar, \.{U}.~Atila, and M.~U\c{c}ar, ``An ensemble approach
for classification of tympanic membrane conditions using soft voting
classifier,'' \textit{Multimedia Tools Appl.}, 2024,
doi:~10.1007/s11042-024-18631-z.

\end{thebibliography}

\end{document}
